% chapters/01/contoh.tex
\section{Contoh Membaca Satu Bab Inti}
\label{sec:contoh-bab}

Sebelum menutup bab ini, ada baiknya Anda melihat bagaimana sebuah bab inti
dirancang, agar Anda dapat membacanya dengan cara yang menguntungkan. Dua
contoh berikut memperlihatkan hal itu.

\subsection{Contoh Pertama: Membaca Keterkaitan Sub-CPMK dengan Asesmen}

Ambil \SUBCPMK{1.2}, yaitu merancang dan mengimplementasikan gerak objek
berbasis \istilah{delta time} dan vektor pada Pygame. Ketika Anda membuka
Bab~V, Anda akan menemukan bahwa Sub-CPMK tersebut disertai tiga indikator
pencapaian, dan setiap indikator berpasangan dengan langkah praktikum serta
bukti yang diperiksa. Pasangan tersebut diperlihatkan pada
Tabel~\ref{tab:contoh-pasangan}.

\begin{table}[htbp]
\centering
\caption{Pasangan indikator, aktivitas, dan bukti pada \SUBCPMK{1.2}}
\label{tab:contoh-pasangan}
{\setlength{\tabcolsep}{4pt}\footnotesize%
\begin{tabularx}{\textwidth}{@{}L{4.6cm} L{4.5cm} Y@{}}
\toprule
\headrow \textbf{Indikator pencapaian} & \textbf{Aktivitas pada praktikum} & \textbf{Bukti yang diperiksa} \\
\midrule
\ind{a} menulis kode gerak menggunakan \code{pygame.Vector2} dan
\istilah{delta time} &
Langkah 1--3: menyusun kerangka program gerak dan menambahkan vektor
kecepatan. &
Berkas \code{code/python/ch05_gerak.py} pada repositori. \\
\addlinespace
\ind{b} menunjukkan kecepatan objek tetap konsisten pada 30 FPS dan 60 FPS &
Langkah 4: menjalankan program pada dua nilai FPS berbeda, lalu mencatat waktu
yang dibutuhkan objek untuk menempuh jarak yang sama. &
Catatan pengamatan pada jurnal praktikum. \\
\addlinespace
\ind{c} menyinkronkan posisi bertipe \istilah{float} dengan \code{Rect}
bertipe bilangan bulat &
Langkah 5: memindahkan hasil perhitungan posisi ke \code{Rect} dan memeriksa
apakah gambar objek bergeser dengan mulus. &
Potongan kode penyelarasan posisi beserta penjelasan singkat. \\
\bottomrule
\end{tabularx}}
\tabsource{Asesmen bab: rubrik portofolio praktikum (kesesuaian fungsional)
dan butir Kuis 1 tentang alasan pemakaian \istilah{delta time}.}
\end{table}

Jika Anda membaca bab dengan memperhatikan tabel semacam ini, tiga manfaat
langsung Anda peroleh. Pertama, Anda mengetahui mengapa sebuah langkah
praktikum harus dikerjakan, bukan sekadar mengikutinya. Kedua, Anda dapat
memeriksa sendiri apakah pekerjaan Anda sudah memenuhi indikator, sebelum
dinilai. Ketiga, ketika nilai Anda belum sesuai harapan, Anda dapat menunjuk
indikator mana yang belum terpenuhi, bukan menebak-nebak.

\subsection{Contoh Kedua: Tuntutan yang Meningkat pada Kemampuan yang Sama}

Kemampuan yang tampak serupa muncul lebih dari sekali dengan
tuntutan yang berbeda. Menggerakkan objek, misalnya, muncul pada Blok I dan
kembali pada Blok III.

\begin{contohbox}
\textbf{Pada Bab~V (Blok I).} Anda diminta menulis sendiri seluruh
perhitungannya: membaca tombol yang ditekan, menyusun vektor arah,
mengalikannya dengan kecepatan dan \istilah{delta time}, lalu menambahkan
hasilnya ke posisi objek. Seluruh keputusan ada di tangan Anda, dan seluruh
kesalahan juga.

\vspace{0.4em}
\textbf{Pada Bab~XIII (Blok III).} Anda diminta mengonfigurasi
\istilah{Rigidbody} dan \istilah{Collider} pada Unity, lalu memeriksa apakah
gerak yang dihasilkan sesuai rancangan. Perhitungannya sudah dikerjakan
\istilah{engine}; yang menjadi tanggung jawab Anda adalah memilih
pengaturan yang tepat dan menganalisis ketika hasilnya menyimpang.
\end{contohbox}

Yang berubah bukan hanya perkakasnya, melainkan juga jenis kesalahan yang
harus Anda temukan. Pada Blok I, kesalahan umumnya berupa perhitungan yang
keliru dan dapat dilacak dari baris kode Anda sendiri. Pada Blok III,
kesalahan umumnya berupa konfigurasi yang tidak sesuai sehingga perilaku objek
terasa janggal---misalnya objek bergetar di permukaan lantai, atau menembus
dinding ketika bergerak cepat. Kemampuan menganalisis gejala semacam itu
dilatih pada \SUBCPMK{3.2} dan menjadi bekal penting pada saat proyek akhir.

Dua contoh di atas memperlihatkan pola yang akan Anda temui berulang kali:
setiap bab menyebut kemampuan yang hendak dicapai, menyediakan latihan untuk
mencapainya, dan memeriksanya dengan cara yang terbuka. Dengan bekal
pemahaman itu, Anda dapat mulai bekerja pada bab berikutnya.
